\naslov{Lastne vrednosti}

Dana je matrika $A \in \R^{n \times n}$.
Iščemo njene lastne vrednosti, in morda lastne vektorje.
Iščemo lahko tudi leve lastne vektorje $y^H A = \lambda y^H$.

\begin{lema}
  Če je $x$ desni lastni vektor $A$ za lastno vrednost $\lambda$ in je $y$ levi
  lastni vektor za lastno vrednost $\mu \ne \lambda$, potem je $y^H x = 0$.
\end{lema}

\begin{proof}
  Velja $\lambda y^H x = y^H Ax = \mu y^H x$, torej $y^H x = 0$.
\end{proof}

\begin{posledica}
  Če je $A$ simetrična in sta $x, y$ lastna vektorja za različni lastni
  vrednosti, je $y^T x = 0$.
\end{posledica}

\vprasanje{Kakšni so lastni vektorji simetrične matrike? Dokaži.}

\begin{vo}{Zakaj Jordanova forma ni primerna za numerično računanje?}
  Poglejmo si primer
  \[
	A(\varepsilon) =
	\begin{bmatrix}
	  0 & 1 \\ \varepsilon & 0
	\end{bmatrix}.
  \]
  Matrika $A(0)$ je že svoja Jordanova kletka, ki pa ni blizu Jordanove forme za
  $\varepsilon \ne 0$, ki je
  \[
	J =
	\begin{bmatrix}
	  \sqrt{\varepsilon} &  \\ & -\sqrt{\varepsilon}
	\end{bmatrix}.
  \]
\end{vo}

\begin{izrek}
  Za vsako matriko $A \in \R^{n \times n}$ obstaja Schurova forma $A = USU^H$,
  kjer je $U$ unitarna in $S$ zgornje trikotna.
\end{izrek}